\scanpage{20}

% loc: top continuation
commutant aux petites $\varinjlim$. Donc $\beta$ induit $\beta':\underline{\operatorname{Hom}}'(E,F)\unsure{-|\to}\underline{\operatorname{Hom}}'(C,F)$. On a pour $\psi$ de $\underline{\operatorname{Hom}}'$

% loc: first calculation
\[\alpha\beta(\psi)(X)=\varinjlim_{Y\in C/X}\psi(Y)\xrightarrow{\sim}\psi(X)\quad\text{car}\quad X=\varinjlim_{Y\in C/X}Y,\ C\]

% loc: faithfulness
étant str. gén. Donc on trouve que $\beta':\underline{\operatorname{Hom}}'(E,F)\to\underline{\operatorname{Hom}}'(C,F)$ est pleinement fidèle. C'est une équivalence i.e. ess. surj. sss $\alpha$ \deleted{\unsure{induit}} envoie bien $\underline{\operatorname{Hom}}'(C,F)$ dans $\underline{\operatorname{Hom}}'(E,F)$. Or soit $(X_i)$ un syst. inductif dans $E$, \deleted{on}\inserted{ayant dans $E$ une limite inductive $Y$} peut alors trouver un système inductif $Y_{ji}$ dans $C$, \inserted{sur l'ens d'indices $I'$ \unsure{est|et}} fibré sur $I$, \unsure{telle|tels|tel} que les \deleted{\unsure{lim}}\inserted{syst. inductifs} partiels (à $i$ fixé) \deleted{\unsure{de $E$}}\inserted{sont} les $C/X_i$, \inserted{de limite $X_i$} et que le système inductif correspondant soit $(Y_{ij})$. On a donc

% loc: iterated colimits
\[\varinjlim_i\deleted{\unsure{\beta}}\alpha(\varphi)(Y_i)=\varinjlim_i\big(\varinjlim_j\varphi(Y_{ij})\big)=\varinjlim_{ij}\varphi(Y_{ij})\underset{\longrightarrow}{\deleted{\unsure{\simeq\lim}}}\varphi(Y)\]

% loc: reduction
et le composé est l'hom dont on veut prouver que c'est un isom. Donc on est ramené à prouver ceci :

% loc: lemma
\uline{Lemme}\quad Si $\varphi\in\underline{\operatorname{Hom}}'(\unsure{C|E},F)$, alors pour $\forall$ système inductif $Y_i$ dans $C$ ayant une limite inductive $Y$ \uline{dans} $E$, on a $\varinjlim\varphi(Y_i)\xrightarrow{\sim}\beta\varphi(Y)$.

% loc: final comparison
Il faut donc comparer \deleted{\unsure{le}} les images par $\varphi$ des 2 systèmes inductifs $(Y_i)$ et $C/Y\to C$. Evidemment on a

% loc: left of cofinality question
\margin{?}

% loc: final question
$I\to C/Y$, cela est-il cofinal ? C'est vrai en tous cas si pour $\forall$ diagramme \uline{fini} dans $C$, la $\varinjlim$ \uline{dans $E$} existe et est dans $C$.

% archive sheet number: 19
